\documentclass[11pt,a4paper]{article}
\usepackage[margin=28mm]{geometry}
\usepackage[T1]{fontenc}
\usepackage{lmodern}
\pdfextension mapfile{+lm.map}
\pdfextension mapfile{+cm.map}
\pdfextension mapfile{+symbols.map}
\pdfextension mapfile{+euler.map}
\usepackage{amsmath,amssymb,amsthm,mathtools}
\usepackage{microtype,longtable,booktabs,array,enumitem}
\usepackage{xcolor}
\usepackage[unicode,colorlinks=true,linkcolor=blue!45!black,citecolor=blue!45!black,urlcolor=blue!45!black]{hyperref}
\usepackage[numbers,sort&compress]{natbib}
\setlength{\emergencystretch}{2em}
\newtheorem{theorem}{Theorem}[section]
\newtheorem{proposition}[theorem]{Proposition}
\newtheorem{lemma}[theorem]{Lemma}
\newtheorem{corollary}[theorem]{Corollary}
\theoremstyle{definition}
\newtheorem{definition}[theorem]{Definition}
\newtheorem{example}[theorem]{Example}
\newtheorem{remark}[theorem]{Remark}
\newenvironment{claim}[3]{\par\medskip\noindent\textbf{#1 (#2). #3.}\ \itshape}{\par\medskip}
\newcommand{\Cat}{\mathrm{Cat}}
\newcommand{\Hot}{\mathrm{Hot}}
\newcommand{\Ho}{\operatorname{Ho}}
\newcommand{\Sp}{\mathrm{Sp}}
\newcommand{\sSet}{\mathrm{sSet}}
\newcommand{\Map}{\operatorname{Map}}
\newcommand{\hocolim}{\operatorname*{hocolim}}
\newcommand{\holim}{\operatorname*{holim}}
\newcommand{\colim}{\operatorname*{colim}}
\newcommand{\cofib}{\operatorname{cofib}}
\newcommand{\fib}{\operatorname{fib}}
\newcommand{\loc}{\operatorname{loc}}
\newcommand{\Fib}{\operatorname{Fib}}
\newcommand{\Catfib}{\mathrm{Catfib}}
\newcommand{\Parf}{\mathrm{Parf}}
\newcommand{\Fun}{\operatorname{Fun}}
\newcommand{\id}{\mathrm{id}}
\newcommand{\RP}{\mathbb{R}P}
\newcommand{\BZ}{B\mathbb Z}
\newcommand{\WZ}{W_{\mathbb Z}}
\newcommand{\Winfty}{W_\infty}
\newcommand{\W}{W}
\newcommand{\cS}{\mathcal S}
\newcommand{\Nul}{\mathfrak N}
\newcommand{\LC}{\mathrm{LC}}
\newcommand{\Grpd}{\mathrm{Grpd}}
\newcommand{\Wh}{\operatorname{Wh}}
\newcommand{\Aut}{\operatorname{Aut}}
\newcommand{\Hom}{\operatorname{Hom}}
\newcommand{\Ob}{\operatorname{Ob}}
\newcommand{\im}{\operatorname{im}}
\newcommand{\coker}{\operatorname{coker}}
\newcommand{\ZZ}{\mathbb Z}
\newcommand{\QQ}{\mathbb Q}
\newcommand{\FF}{\mathbb F}
\title{Reversible localizers: proofs, counterexamples, and an audit of the proposed constructions}
\author{AI-produced by GPT-6 Astra\\OpenAI, ChatGPT Pro mode}
\date{15 November 2026}
\begin{document}
\maketitle
\noindent\textbf{Working status: incomplete.} We claim no priority. Each numbered statement below carries its own status. Unresolved requested assertions are kept distinct from proved reductions.

\paragraph{Conventions.}
We fix a universe $\mathcal U$ and let $\Cat$ denote its small categories. We put $BX=|NX|$, and $\Winfty$ consists of functors inducing weak homotopy equivalences on nerves. Unless Part F says otherwise, a fundamental localizer means Cisinski's definition: weak saturation, the terminal-object axiom, and relative Theorem A. The latter says that for a triangle $A\xrightarrow{u}B\to C$, if $A/c\to B/c$ belongs to $W$ for every $c\in C$, then $u\in W$. A category $X$ is $W$-aspherical when $X\to e$ belongs to $W$. A sieve is closed under sources, and $A/c$ uses arrows from the image of $A$ to $c$. Homotopy cofibres of unpointed maps carry their cone basepoint; suspension spectra of pointed spaces are reduced. All mapping spaces and homotopy limits and colimits are derived. Statements include empty and disconnected categories unless an explicit hypothesis excludes them.
\tableofcontents
\InputIfFileExists{sections/A.tex}{}{}
\InputIfFileExists{sections/B.tex}{}{}
\InputIfFileExists{sections/C.tex}{}{}
\InputIfFileExists{sections/D.tex}{}{}
\InputIfFileExists{sections/E.tex}{}{}
\InputIfFileExists{sections/F12.tex}{}{}
\InputIfFileExists{sections/F3.tex}{}{}
\InputIfFileExists{sections/G.tex}{}{}
\InputIfFileExists{sections/H14.tex}{}{}
\InputIfFileExists{sections/H5.tex}{}{}
\InputIfFileExists{sections/I.tex}{}{}
\InputIfFileExists{sections/J.tex}{}{}
\bibliographystyle{plainnat}
\bibliography{REVLOC_references}
\end{document}
